\documentclass[11pt,a4paper]{article}
\usepackage[margin=2.2cm]{geometry}
\usepackage[T1]{fontenc}
\usepackage{lmodern}
\usepackage{amsmath}
\usepackage{xcolor}
\usepackage{tabularx}
\usepackage{booktabs}
\usepackage{enumitem}
\usepackage[hidelinks]{hyperref}
\usepackage{titlesec}
\definecolor{navy}{HTML}{0E2841}
\definecolor{teal}{HTML}{156082}
\definecolor{tint}{HTML}{E6EFF5}
\titleformat{\section}{\Large\bfseries\color{navy}}{\thesection}{0.6em}{}
\titleformat{\subsection}{\large\bfseries\color{teal}}{\thesubsection}{0.5em}{}
\setlength{\parskip}{4pt}
\setlength{\parindent}{0pt}
\setlist{itemsep=2pt,topsep=3pt}
\newcommand{\slide}[3]{\subsection*{Slide #1: #2 \hfill\normalfont\small\color{teal}about #3}}
\newcommand{\say}[1]{\par\noindent\colorbox{tint}{\parbox{\dimexpr\linewidth-2\fboxsep}{#1}}\par}
\newcommand{\q}[1]{\par\medskip\noindent\textbf{\color{navy}Q: #1}\par}
\newcommand{\ans}{\noindent\textbf{A:} }

\begin{document}

\begin{center}
{\LARGE\bfseries\color{navy} Presentation Guide}\\[4pt]
{\large Parameter-Efficient UI Code Generation Using QLoRA on an Extended WebSight Dataset}\\[4pt]
ISASF-2026 \quad|\quad Paper ID 132 \quad|\quad 19 November 2026 (Hybrid)\\[2pt]
{\small Aniruddha Dev Nath, Anitya Barua, Antu Chowdhury, Arnab Biswas}
\end{center}

\medskip
\noindent\textbf{How to use this guide.} Part~1 is the 5-minute speech, slide by slide;
the shaded text is what to say. Part~2 explains every metric and term. Part~3 lists the questions
reviewers and the audience are most likely to ask, with honest answers. Part~4 covers general questions.
Every number here is taken from the final camera-ready paper.

\medskip
\noindent\textbf{Three rules for the talk.}
\begin{enumerate}
\item Say ``statistically tied'' for SSIM. The fine-tuned 2B model is \emph{not} better than the 3B models on SSIM.
\item Say ``saves memory, not time''. The 4-bit model is slower than the 16-bit base model.
\item On real pages, the fine-tuned model writes much shorter code than the original pages. Say this yourself before someone asks.
\end{enumerate}

%======================================================================
\section{Speech: 5-minute version (9 slides)}

Use the deck \texttt{ISASF2026\_Paper132\_Presentation\_5min.pptx}. Speak slowly: the script is about
4 minutes 15 seconds, which leaves time for changing slides. If you run short of time, skip the sentences marked
\textit{(optional)}.

\slide{1}{Title}{15 s}
\say{Good morning. I am [name]. Our paper asks one question: can a small vision-language model, fine-tuned on a
single 8-gigabyte GPU, turn a screenshot of a user interface into React or Vue code?}

\slide{2}{Introduction}{40 s}
\say{Screenshot-to-code research mostly produces HTML, but real frontends are built from React and Vue components,
and large commercial models cost money per token and see your unreleased designs.
No dataset pairs screenshots with validated React and Vue code. And small models struggle here: used zero-shot,
the 2B model imports UI libraries in 437 of its 496 React outputs, even though we asked for Tailwind only.}

\slide{3}{Methodology}{50 s}
\say{We did three things. First, we extended WebSight: a code model generated React and Vue versions of every page,
each checked by a real parser; 15,588 samples have both targets valid.
Second, we fine-tuned two models with QLoRA on one 8-gigabyte GPU: Qwen2-VL-2B, and for comparison Qwen2.5-VL-3B.
Third, we evaluated the code, and we rendered every output in a browser to compare it with the original screenshot,
with bootstrap confidence intervals.}

\slide{4}{Results: WebSight}{55 s}
\say{On 1,488 test samples, fine-tuning the 2B model raises BLEU-4 from 52 to 81 and CodeBLEU from 0.52 to 0.74.
Against the larger zero-shot 3B model, our renders are significantly closer on LPIPS and PSNR.
On SSIM the two are statistically tied; SSIM is near its ceiling of 0.722.
The fine-tuned 3B improves a lot too, but does not beat the 2B. \textit{(optional:) It had only one epoch instead of three.}}

\slide{5}{Efficiency}{35 s}
\say{Our 2B model needs 2.18 gigabytes of GPU memory, a third of the zero-shot 3B model, and it is about twice as fast.
But to be clear: 4-bit QLoRA saves memory, not time. It is 1.75 times slower than the 16-bit base model.}

\slide{6}{Real web pages (Design2Code)}{45 s}
\say{Reviewers asked about real websites, so we tested 100 real pages from Design2Code.
Fine-tuning still helps strongly: SSIM rises by 0.21 over the base model.
But here our model is only tied with the 3B models, and it writes about 600 tokens per page, while real HTML has a
median of nearly 15,000. It learned the length of the short synthetic training pages.}

\slide{7}{Conclusion}{30 s}
\say{To conclude: we built more than 15,000 parser-validated React and Vue targets, and showed that a 2B model
fine-tuned on one consumer GPU can match a larger fine-tuned model and, on WebSight-style pages, beat a larger zero-shot
one, with designs kept on-premises.
\textit{(optional:) Next, we want to train on longer real pages and compare QLoRA with LoRA and DoRA.}}

\slide{8--9}{References and Thank You}{5 s}
\say{Thank you. I am happy to take your questions.}

\medskip
\noindent\textbf{Timing check:} 15 + 40 + 50 + 55 + 35 + 45 + 30 + 5 seconds = 4 minutes 35 seconds.
The longer 14-slide deck (\texttt{ISASF2026\_Paper132\_Presentation.pptx}) stays as a backup; the slides removed from
the 5-minute version (abstract, literature review, dataset, discussion, limitations) are useful if someone asks.

%======================================================================
\section{What every metric and term means}

\subsection*{Code-quality metrics (compare the generated code with the reference code)}
\renewcommand{\arraystretch}{1.25}
\begin{tabularx}{\linewidth}{@{}p{3.1cm}X p{2.3cm}@{}}
\toprule
\textbf{Metric} & \textbf{Meaning} & \textbf{Better} \\
\midrule
Parse Success Rate (PSR) & Percentage of outputs that a real parser accepts: \texttt{@babel/parser} for React JSX,
\texttt{@vue/compiler-sfc} for Vue. It measures whether the code is syntactically valid. & Higher (0--100\%) \\
BLEU-4 & Counts how many sequences of 1 to 4 tokens in the output also appear in the reference code. Originally a
machine-translation metric. & Higher (0--100) \\
CodeBLEU & BLEU adapted for code: it also compares the syntax tree (AST) and the data flow of the code. Computed with a
fixed Python hash seed, because its data-flow part depends on it. & Higher (0--1) \\
Edit similarity & $1 - \frac{\text{edit distance}}{\text{length of the longer text}}$: how many character edits are
needed to turn the output into the reference. & Higher (0--1) \\
Tag accuracy & Fraction of the reference HTML/JSX tags that also appear in the output. & Higher (0--1) \\
Tree edit similarity & Edit similarity computed on the sequence of markup tags instead of characters: compares the
page structure. & Higher (0--1) \\
\bottomrule
\end{tabularx}

\subsection*{Visual metrics (render the code in a browser, compare with the original screenshot)}
\begin{tabularx}{\linewidth}{@{}p{3.1cm}X p{2.3cm}@{}}
\toprule
\textbf{Metric} & \textbf{Meaning} & \textbf{Better} \\
\midrule
SSIM (structural similarity) & Compares local brightness, contrast and structure of two images. 1 means identical.
Our upper bound (the reference code itself) is 0.722. & Higher (0--1) \\
PSNR (peak signal-to-noise ratio) & $10\log_{10}(\mathrm{MAX}^2/\mathrm{MSE})$ in decibels; based on the pixel-by-pixel
error. & Higher (dB) \\
LPIPS (learned perceptual similarity) & Distance between deep-network (AlexNet) features of the two images; closer to
how humans judge similarity. & \textbf{Lower} (0--1) \\
Rendered (\%) & Percentage of outputs that produce a visible page in the browser. & Higher \\
\bottomrule
\end{tabularx}

\noindent Failed renders are \emph{not} removed: they get the worst score (SSIM 0, PSNR 0, LPIPS 1), so visual scores
reflect both renderability and appearance. Renders are made at a 1280$\times$960 viewport in headless Chromium.

\subsection*{Statistics and agreement}
\begin{tabularx}{\linewidth}{@{}p{3.1cm}X@{}}
\toprule
Paired bootstrap & We resample the 496 test images with replacement (1,000 times; 500 for CodeBLEU) and recompute
the difference between two models each time. ``Paired'' means both models are compared on the same resamples, and
the three tasks of an image are resampled together. \\
95\% confidence interval (CI) & The range that contains the true difference with 95\% confidence. If it does not
include 0, the difference is significant. \\
$p$-value & Two-sided. $p<0.05$ means significant. Example: SSIM 2B vs.\ zero-shot 3B has $p=0.33$, so it is a tie. \\
PABAK & Prevalence-Adjusted Bias-Adjusted Kappa: agreement between validators, corrected for the fact that almost all
samples are valid (where Cohen's kappa is misleading). Above 0.81 is ``almost perfect''; ours is 0.9786. \\
\bottomrule
\end{tabularx}

\subsection*{Model and training terms}
\begin{tabularx}{\linewidth}{@{}p{3.1cm}X@{}}
\toprule
VLM & Vision-language model: takes an image and text, produces text. \\
Zero-shot & The model is used as released, without any fine-tuning on our data. \\
LoRA & Low-Rank Adaptation: the original weights are frozen; only two small matrices (rank $r$) are trained per layer. \\
QLoRA & LoRA where the frozen weights are stored in 4-bit (NF4) to save memory. We used rank $r=64$, $\alpha=128$. \\
NF4 & 4-bit NormalFloat, the 4-bit number format used by QLoRA. \\
FP16 / BF16 & 16-bit number formats. The 2B adapters were trained in FP16; the 3B in BF16, because FP16 overflowed. \\
Epoch & One full pass over the training data. 2B: 3 epochs; 3B: 1 epoch (3,152 steps, about 34 hours). \\
Teacher model & Qwen2.5-Coder-7B, which generated our React and Vue training targets. \\
Self-repair loop & If a generated target fails the parser, the error message is sent back to the teacher, which
regenerates it (up to two attempts). \\
Max new tokens & Generation limit of 1,536 tokens; it covers 99.93\% of the test references. \\
Design2Code & A benchmark of 484 real web pages. We used a random 100, first screen only (16:9 crop). \\
Lenient harness & A rendering variant that replaces imported UI-library components with plain \texttt{<div>}
placeholders, to separate broken code from missing libraries. \\
\bottomrule
\end{tabularx}

%======================================================================
\section{Likely questions from reviewers and the audience}

\q{Your training targets come from another model. Isn't the evaluation circular?}
\ans Partly, for the code metrics, and we say so. That is why we treat the visual metrics as primary. Four results
limit the concern: (1) the HTML references are the original WebSight code, not from our teacher, and fine-tuning still
raises HTML BLEU-4 from 64.69 to 83.23; (2) visual metrics compare with the original screenshot, not the teacher's code,
and our model reaches 99\% of the reference code's SSIM; (3) a render check shows the teacher's targets reproduce their
source pages; (4) on Design2Code, where no teacher code exists, fine-tuning still improves every visual metric.

\q{Why is your model not better than the 3B model on SSIM?}
\ans SSIM is statistically tied ($p=0.33$ against the zero-shot 3B, $p=0.87$ against the fine-tuned 3B). SSIM is close
to its ceiling: even the reference code only reaches 0.722, so it separates good models poorly. On LPIPS and PSNR our 2B
model is significantly better on WebSight.

\q{Why did the fine-tuned 3B model not beat the 2B model?}
\ans It was trained for one epoch instead of three, and in BF16 instead of FP16, because of time and memory limits. Its
training loss at the same steps was lower than the 2B's, so it may simply be under-trained. We did not test this, so we
do not claim that the larger model is worse in general.

\q{QLoRA is supposed to be efficient. Why is it slower than the base model?}
\ans QLoRA saves memory, not time. The 4-bit weights must be dequantized at every step, which is a known overhead for
QLoRA models, and the adapter is applied separately. At the same batch size, our model needs 11.2\,s per sample versus
6.4\,s for the 16-bit base model, but uses 2.18 versus 4.59\,GiB of memory. Earlier latency numbers that looked faster
came from different batch sizes; we re-measured under identical conditions.

\q{Does it work on real websites?}
\ans Partly. On 100 real Design2Code pages, fine-tuning raises SSIM by 0.210 over the base model ($p<0.001$), and the
React render rate rises from 13\% to 93\%. But all scores drop, our model is tied with both 3B models there, and it
writes about 600 tokens per page while the real HTML has a median of 14,805 tokens. Training on longer real pages is our
main future work.

\q{Why does the base model render only 12.1\% of its React outputs?}
\ans Mostly because it imports UI libraries such as react-bootstrap (437 of 496 outputs) despite the instruction to use
only Tailwind CSS, and our harness does not install those libraries. With a lenient harness that replaces them with
placeholders, 87.7\% render and the SSIM gap to our model shrinks to 0.027. So much of the strict gap is about following
instructions, and we report both settings.

\q{Could the test set contain near-copies of the training pages?}
\ans We checked. No test page duplicates a training page; one exceeds the common near-duplicate threshold of 0.8 (word
5-gram Jaccard similarity) and 28 exceed 0.5. Removing these 28 pages changes SSIM, LPIPS and CodeBLEU by less than 0.01
and does not change the ranking of the models.

\q{Why did you not compare with GPT-4o or other commercial models?}
\ans Our goal was a model that runs on a single consumer GPU without per-token cost and without sending designs to a
third party, so we compared against open models of similar size. We did not test commercial models, so we make no claim
about how we compare with them.

\q{Why no few-shot baseline, or an existing WebSight-trained model?}
\ans A few-shot prompt needs at least one example screenshot plus its full code (about 1,500 tokens per example). The
zero-shot 3B model in 16-bit already needs more memory than our 8\,GB GPU provides. The WebSight-trained models we know
of have 7B parameters or more and generate HTML only, so they cannot run on our hardware or produce React and Vue.

\q{Did you run several seeds? How reliable are the results?}
\ans Each model was trained once, which is a limitation. To quantify uncertainty we report paired bootstrap 95\%
confidence intervals and $p$-values over the 496 test images for every comparison.

\q{How do you know the generated React and Vue training targets are correct?}
\ans Every target passed a production parser, with up to two repair attempts. Three independent validators agree with a
mean PABAK of 0.9786. We also rendered the targets of 2,000 images: 99.5\% of the JSX and 100\% of the Vue targets render,
with a median SSIM of 0.90 and 0.96 to their source page.

\q{Is the generated code ready for production?}
\ans No. It handles static layouts, not state management, interaction or accessibility, and it should be reviewed by a
developer before it is merged.

\q{Why did you limit generation to 1,536 tokens? Does that cut outputs?}
\ans The limit covers 99.93\% of the test references. None of our fine-tuned model's WebSight outputs reached it; it cut
3.7\% of the base model's and 1.1\% of the zero-shot 3B model's outputs.

%======================================================================
\section{General questions}

\q{Why React and Vue?}
\ans They are among the most widely used component-based frontend frameworks. HTML output must be converted by hand
before it can be used in such a codebase.

\q{Why Qwen2-VL-2B?}
\ans It is a small open vision-language model that we could fine-tune with QLoRA on an 8\,GB GPU. Qwen2.5-VL-3B is a
newer and larger model from the same family, which lets us compare scale.

\q{What hardware did you use?}
\ans One NVIDIA GeForce RTX 3070 Ti with 8\,GB for all training and inference. The 3B model's single epoch took about 34
hours.

\q{Why WebSight and not real web pages for training?}
\ans WebSight provides clean, self-contained HTML pages with matching screenshots, without the broken layouts and
copyright problems of scraped pages. Its limitation, shorter and simpler pages, shows up in our Design2Code results.

\q{How large is the dataset?}
\ans We processed 15,752 WebSight samples; 15,588 have valid React and Vue targets. For fine-tuning we used 9,889 samples,
giving 25,215 training, 2,964 validation and 1,488 test examples (each image has an HTML, a JSX and a Vue task).

\q{What would you do next?}
\ans Compare QLoRA with LoRA and DoRA on the same hardware, train on longer real-world pages such as WebCode2M, add
interaction-aware metrics, and add a self-repair loop driven by validation errors at inference time.

\q{If you do not know the answer to a question}
\ans Say: ``That is a good point; we did not test that in this paper, but it is a natural next step.'' Do not guess
numbers.

\end{document}
